% SPDX-License-Identifier: CC-BY-4.0
% Copyright (c) 2026 Martin Schröder <info@swedishembedded.com>

\section{Audit of the existing stack against the protocol}\label{sec:audit}

Before building, we read the three code bases against the requirements of
\cref{sec:protocol}. The list below is what that reading found. These are
\emph{static findings}: each was located in source and none has yet been reproduced by a
failing test; an entry moves to a verified defect only when its regression test exists, and
the section is revised then.

\begin{enumerate}
\item \textbf{Loss masking of tool conversations.} The trainer's per-record masking
checks that each supervised span lies on a boundary of the rendered chat. Under the
Qwen3 template, an empty reasoning block is added to the last assistant turn only, so
non-final assistant turns of a multi-step conversation fail the boundary check; turns
with tool calls take that same path. Training on whole agent trajectories therefore needs
a per-answer prompt/completion encoding, with each completion equal to what the serving
renderer produces.
\item \textbf{Tool definitions lost between recording and training.} The agent records
the tool definitions of the run, but the trajectory-to-training projections we inspected
drop them, so a model would be trained on tool calls without the tool block it sees when
served. Splinter's projection also refuses a second user turn (the follow-up turns of the
solver) and substitutes its own system prompt.
\item \textbf{Trajectory fidelity in the agent SDK.} The exported trajectory has no marker
for a context compaction, lists every registered tool instead of the set offered in the
agent's mode, and turns unparseable tool-call arguments into an empty object. Each
changes what a training record would claim the model saw or did.
\item \textbf{Closed-book release probes.} The release gate's probes and serve check run
tasks closed-book; a suite of workspace tasks would be unmeasured and so fail the gate,
and text agreement between served and in-process answers is meaningless across
multi-step runs.
\item \textbf{Containment.} A coding agent's shell is not a sandbox; the file tools are
confined to a root but a command can read outside it. In the environment of this study the
container runtime and the user-namespace sandbox are both unusable, so runs have
isolation level \code{none}, and the exposure audit is detection, not prevention. A
previous coding loop also recorded a case in which a worker read a hidden test.
\item \textbf{Contamination paths.} A semantic-memory tool was not in the worker's
withheld set, and the agent runtime writes its audit log into the workspace.
\item \textbf{Serving.} Brain serves one adapter per process and selects it by manifest
digest; the manifest has a reader but no public writer.
\end{enumerate}
